\section{What the published comparison can and cannot settle}
\label{sec:cannot}



The published comparison ran four arms over the same frozen pipeline, input and recorded
corpus, differing in what the reviewer sees: stage metadata, accumulated history, the
complete graph, or a selected part of it (table in Appendix~\ref{app:published}). It is
carefully controlled: the same frozen pipeline, input
and recorded corpus across arms, fresh sessions, no cross-arm visibility. Three properties
nonetheless prevent it settling which regime is better, and they determine what this paper
had to build.

Sample size is the first. One run per arm and three repairs within each makes every
quantity a count out of three or a single ordinal, and the recommended arm ties the
accumulated-history baseline on effective repairs. Second, the arms were not adjudicated
identically: the archived record states the first common full-graph judge exceeded the
model's 1{,}048{,}576 character input limit on the history branch, and those runs were
rejudged with isolated graph-selected judges, a failure that also removed that arm's
timing. Third, and most usefully, the function selecting which boundary to repair differs
across arms. The graph arms filter flagged units to causal roots using upstream
identifiers; the baselines rotate with no causal step. The comparison is presented as one
of context regimes but the arms also differ in how they choose a target. Rather than
control that away we make it the treatment, and it produces the paper's largest effect.
